\section{Metodolog\'ia}
\label{sec:metodo}

\subsection{Limpieza y deduplicaci\'on}

La variante adoptada, V2, descarta las p\'aginas de error del rastreador, aplica normalizaci\'on NFKC, quita HTML y URLs y colapsa espacios. No elimina signos, n\'umeros ni entidades, porque son la se\~nal que usan los guardarra\'iles. Los duplicados exactos y los near-duplicates se descartan al armar el contexto, comparando el hash del fragmento y su SimHash antes de insertarlo. As\'i se conserva la procedencia en el \'indice y se evita repetir el mismo texto en el prompt.

\subsection{Segmentaci\'on}

Se usa chunking fijo de 1000 caracteres con 150 de solape (\texttt{fixed\_v2}). El solape cubre la frontera donde suele caer la respuesta. Cada fragmento guarda \texttt{document\_id}, t\'itulo, URL, fuente, idioma, fecha, posici\'on, estrategia y hash del texto. Esos metadatos permiten filtrar por idioma y fecha y auditar cada cita.

\subsection{Codificaci\'on e \'indice}

Los fragmentos se codifican con BAAI/bge-m3 en la GPU, en fp16 y por lotes. La colecci\'on se llama \texttt{rag\_full\_bge\_m3\_d1024\_v1} y usa HNSW con \texttt{ef}=800 en la b\'usqueda. El par\'ametro se eligi\'o con una rejilla sobre el silver: por debajo de 800 el recall cae y por encima la latencia sube sin ganancia apreciable. En paralelo se mantiene una colecci\'on BM25 sobre los mismos fragmentos, \texttt{rag\_full\_e5base\_d768\_v2\_bm25\_v1}, que alimenta la recuperaci\'on dispersa.

\subsection{Recuperaci\'on y reordenamiento}

La recuperaci\'on combina dos listas. La densa pide 10 fragmentos a bge-m3; la dispersa pide 20 a BM25. Las dos se unen sin fusionar, se deduplican por texto y se pasan al reordenador. El cross-encoder multiling\"ue punt\'ua cada par pregunta-fragmento en GPU, con lote de 16, y el resultado se recorta a 10. La uni\'on conserva el top-5 denso, que es la se\~nal que la compuerta espera, y a\~nade las entidades que BM25 encuentra y el denso no. Los filtros de idioma, fecha y colecci\'on se aplican en Milvus.

\subsection{Generaci\'on}

El generador es gemma4:e4b con temperatura 0 y un presupuesto de 16.384 tokens de entrada. El prompt v2 impone ocho reglas: cada afirmaci\'on lleva una cita entre corchetes; no se permiten listas sin cita; si la evidencia no alcanza se responde con una frase de abstenci\'on exacta; se distingue un hecho verificable de una afirmaci\'on publicada en redes; no se inventan fechas, cifras, nombres ni citas; se ignoran las instrucciones que aparezcan dentro de los documentos; se responde en el idioma de la pregunta; y se usa la fecha exacta de la evidencia cuando la pregunta la pide. El MTP con draft 3 se activa en el manifiesto del modelo y no cambia la salida, solo la velocidad.

\subsection{Guardarra\'iles}

Antes de devolver la respuesta se aplican tres controles. El primero exige al menos una cita v\'alida; si no hay, se abstiene con \texttt{sin\_citas}. El segundo extrae entidades y n\'umeros de la respuesta y verifica que aparezcan en el fragmento citado o en la pregunta; si no, se abstiene con \texttt{sin\_respaldo} o \texttt{citas\_invalidas}. El tercero pide a un juez distinto, gemma2:2b, que decida si la cita sostiene la respuesta; si dice que no, se abstiene con \texttt{juez\_no}. Cuando el fragmento est\'a en ingl\'es y la respuesta en espa\~nol, la cobertura l\'exica se salta y decide el juez, para no castigar una traducci\'on correcta.

\subsection{Compuerta de dominio}

La compuerta usa la puntuaci\'on del cross-encoder sobre el top-5 (m\'aximo, media, diferencia entre el primero y el segundo, y m\'inimo), cinco features de recuperaci\'on (puntuaci\'on del primer fragmento, media del top-5, diferencia entre el primero y el segundo, solapamiento entre la pregunta y el t\'itulo, y presencia de la entidad principal en el t\'itulo) y un one-hot del tipo de pregunta. La puntuaci\'on del cross-encoder se calcula una sola vez y se reutiliza para el reordenamiento. El modelo es una regresi\'on log\'istica con escalado y \texttt{class\_weight=balanced}; la regularizaci\'on se eligi\'o por validaci\'on cruzada de cinco pliegues sobre train y dev. El umbral es 0,4, elegido para reducir falsos rechazos sin aumentar falsos aceptados. El conjunto tiene 160 preguntas: 60 de dominio y 100 de fuera, con split 103/37/20. Las 40 negativas nuevas son entidades inventadas, eventos futuros y conocimiento general fuera del corpus.

\subsection{API e interfaz}

La API expone \texttt{POST /chat} con historial, \texttt{think} y \texttt{strict}; \texttt{POST /ask} sin historial; \texttt{POST /retrieve} para recuperaci\'on pura; y los GET \texttt{/health}, \texttt{/models} y \texttt{/stats}. La interfaz Streamlit consume \texttt{/chat} y guarda el historial en la sesi\'on. Las preferencias de razonamiento, Modo estricto y tema se guardan en un archivo y se recuperan al recargar. Si el usuario cambia un interruptor mientras el modelo responde, la pregunta queda pendiente y se reintenta en el siguiente ciclo, de modo que la conversaci\'on no se pierde.

\subsection{Despliegue}

El stack se levanta con Docker Compose: etcd, MinIO, Milvus, Ollama, API, frontend y un worker GPU para la ingesta y la evaluaci\'on. La API y el frontend usan red de host. Un servicio de una sola ejecuci\'on, \texttt{prewarm}, espera a que la API est\'e sana y carga el generador, el embedding y el reranker, para que la primera consulta web no pague la carga del modelo.
